\paragraph{A fixed qualitative example.}
Figure~\ref{fig:qualitative-goop2d} shows the metadata-selected source 12 at
forecasts 1, 200 and 395. The ground-truth particles settle near the lower
boundary, while both mixed-training forecasts retain elevated groups at the
final forecast. Thus fewer visible excursions do not establish physically
accurate dynamics. This single example illustrates the remaining mismatch;
it does not estimate population performance. Comparing the base-trained/base
column with either mixed-training column changes both training and evaluation
graphs. The training-effect inference instead rests on the paired comparison
under the same random25 evaluation policy. Figure counts use any strict
crossing of the metadata box, whereas Table~\ref{tab:goop-all-cost} uses
excursions greater than $10^{-6}$. Boundary crosses are clamped markers and do
not show excursion magnitude.
